\documentclass[11pt]{article}

\usepackage[spanish]{babel}
\usepackage[margin=1in]{geometry}
\usepackage{amsmath, amssymb, amsfonts}
\usepackage{booktabs}
\usepackage{enumitem}
\usepackage{tcolorbox}
\usepackage{hyperref}
\usepackage{listings}
\usepackage{xcolor}
\usepackage{parskip}
\usepackage{pgfplots}

\usetikzlibrary{arrows.meta}
\pgfplotsset{compat=1.18}

\lstset{
  language=Python,
  basicstyle=\ttfamily\footnotesize,
  breaklines=true,
  columns=flexible,
  keepspaces=true,
  commentstyle=\color{gray}\itshape,
  keywordstyle=\color{blue!70!black}\bfseries,
  stringstyle=\color{green!40!black},
  frame=L,
  xleftmargin=2em,
  showstringspaces=false,
  tabsize=4
}

\title{\textbf{Hoja de Trabajo 1: Redes Neuronales} \vspace{1em}\\
  Deep Learning y Sistemas Inteligentes - CC3092}
\author{José Antonio Mérida Castejón}
\date{\today}

\begin{document}

\maketitle

\section*{Task 1}
\textit{Responda de forma clara:}

\begin{enumerate}

  \item \textit{Explique con sus propias palabras por que una red neuronal compuesta únicamente de transformaciones lineales, sin funciones de
    activación no lineales, tiene exactamente la misma capacidad expresiva que un modelo lineal de una sola capa. Su explicación debe incluir
    el argumento algebráico central.}

  En términos resumidos, la composición de capas lineales resulta en otra capa lineal. Esto quiere decir que no importa cuantas veces multipliquemos
  la entrada por una matriz de pesos, el resultado seguirá siendo el mismo que multiplicarlo por una única matriz. Además, sin importar cuantos vectores
  de \textit{bias} sumemos seguiremos obteniendo otro vector de \textit{bias}. Para una demostración más formal consideremos lo siguiente:

  Una capa lineal aplica a la entrada $x \in \mathbb{R}^n$ una transformación de la forma
  \[
    f_1(x) = W_1 x + b_1,
  \]
  donde $W_1 \in \mathbb{R}^{k \times n}$ es la matriz de pesos y $b_1 \in \mathbb{R}^k$ el \textit{bias}, de modo que $f_1$ va de $\mathbb{R}^n \to \mathbb{R}^k$.

  Ahora, consideremos una segunda capa $f_2: \mathbb{R}^k \to \mathbb{R}^m$, con $W_2 \in \mathbb{R}^{m \times k}$ y $b_2 \in \mathbb{R}^m$:
  \[
    f_2(y) = W_2 y + b_2.
  \]
  Al componer ambas sin una función de activación de por medio:
  \[
    f_2(f_1(x)) = W_2 (W_1 x + b_1) + b_2 = (W_2 W_1)\, x + (W_2 b_1 + b_2).
  \]
  El producto $W_2 W_1$ es de nuevo una matriz $W' = W_2 W_1 \in \mathbb{R}^{m \times n}$. De igual forma, $b_1$ es multiplicado por $W_2$,
  que lo lleva de $\mathbb{R}^k$ a $\mathbb{R}^m$, y al sumarle $b_2$ obtenemos otro vector $W_2 b_1 + b_2 = b' \in \mathbb{R}^m$. Por lo tanto:
  \[
    f_2(f_1(x)) = W' x + b',
  \]
  que tiene la misma forma que una única capa $f': \mathbb{R}^n \to \mathbb{R}^m$. Repitiendo el argumento, cualquier número de capas colapsa en una sola $W'x + b'$, por lo que la red tiene la misma capacidad expresiva que un modelo lineal de una sola capa.

  \item \textit{Considere una neurona con dos entradas $x_1 = 3$ y $x_2 = -2$, pesos $w_1 = 0.4$ y $w_2 = 0.6$, y sesgo $b = -0.1$. Calcule la
    preactivación $zy$ y la activación $a = \sigma(z)$. Interprete el resultado, ¿Qué le indica este valor sobre la predicción del modelo?}
  \[
    z = w_1 x_1 + w_2 x_2 + b = (0.4)(3) + (0.6)(-2) + (-0.1) = -0.1.
  \]
  \[
    a = \sigma(z) = \frac{1}{1 + e^{-z}} = \frac{1}{1 + e^{0.1}} \approx 0.4750.
  \]

  La activación $a \approx 0.475$ se interpreta como la probabilidad estimada de que la entrada pertenezca a la clase positiva, es decir,
  $P(y = 1 \mid x) \approx 0.475$. Como este valor está apenas por debajo de $0.5$, el modelo predice la clase negativa (\textit{no}), pero por un
  margen muy pequeño.

\end{enumerate}


\end{document}
